% GENERATED, edit the French unit or the translation table
\mxfront{premiere-front}
\begin{titlepage}
\color{mxink}\sffamily
{\large MATHS EXPLORER\par}
\vspace{22mm}
{\fontsize{35}{40}\selectfont\bfseries Mathematics\par}
\vspace{6mm}
{\LARGE Première\par specialism\par}
\vspace{12mm}
{\Large 2026--2027\par}
\vspace{10mm}
\begin{tikzpicture}[scale=.8]\draw[mxmint,very thick,domain=-1.7:1.7,samples=80]plot(\x,{\x*\x});\draw[mxink,thick](-2,0)--(2,0);\draw[mxink,thick](0,-.3)--(0,3.2);\draw[mxwarn,thick](-.5,-2)--(2,3);\end{tikzpicture}
\vfill
{\Large\bfseries\ifmxteacher Teacher edition\else Student edition\fi\par}
\vspace{4mm}
Developed lessons, progressive pathways and preparation for the early exam.\par
Ten chapters in the order of the school's progression.\par
\vspace{8mm}
{\small Editorial version revised on 10 September 2026.\par To be used with CQFD Mathématiques 1re, Bordas, 2019.}
\end{titlepage}
\chapter*{Using this booklet}
\addcontentsline{toc}{chapter}{Using this booklet}
The lessons contain fully worked examples and the proofs to know. The diagnostics, applications, exercises, MCQs and practice sets are to be done in a notebook. The teacher has the solutions and reserved variants. The booklet follows the progression supplied; the textbook numbers are kept in the references and must not be confused with the numbers of this booklet.

Each chapter begins with six prerequisite questions and an investigation activity. The pathway then leads from guided exercises to the choice of a method, then to transfer. Exercise 8 is intended for a delayed review. The reserve exercises are not a list of homework to be completed in full.

\begin{keypoint}[How to work]
Investigate first without a solution. Write down the data, the domain and the question. After help or a correction, close the lesson and rebuild one step. The following week, redo a short task without notes. A wrong answer accompanied by a usable attempt lets the teacher target the help.
\end{keypoint}

\textbf{Writing up.} An equality is justified by a definition, a property or a calculation. A conclusion answers the question, with a unit if necessary. ``Conjecture'' means to suggest a plausible answer; ``prove'' means to establish its validity under explicit hypotheses.

\textbf{Tools.} Numerical exercises may use Python when asked. For the fluency exercises and the EAM preparation papers, calculators are not allowed. For homework, keep the first attempt and note any help received; being able to explain one's solution without tools is a central goal.

\ifmxteacher\paragraph{Feedback on assessments.} Keep three indicators: accuracy, choice of method, justification. A correct MCQ does not certify the write-up. Redo an incorrect task without the lesson, then after a week and in the consolidation slot.\fi

\textbf{Notation.} Intervals follow the French conventions. A conditional probability may be written $P_A(B)$ or $P(B\mid A)$. Geometric figures are in an orthonormal frame when the calculations require it. The figures in the lessons illustrate a property; they do not replace its proof.

\chapter*{The year and the EAM framework}
\addcontentsline{toc}{chapter}{The year and the EAM framework}
The programme used is that of the order of 26 February 2026 (NOR MENE2602917A), published in BO no. 14 of 2 April 2026. Its article 2 makes it applicable from the start of the 2026--2027 school year. It replaces the annex of the 2019 Première programme. The new Terminale programme planned for the 2027--2028 school year is not the one used in this booklet.

Première students in 2026--2027 sit the early exam (EAM) in 2027, for the 2028 baccalauréat. The ministerial note of 10 June 2025 (NOR MENE2515469N) defines a two-hour exam, weighting 2, without a calculator: a fluency MCQ worth 6 points, then two or three independent exercises worth 14 points. The transitional list of fluency skills explicitly concerning 2025--2026 does not replace the fluency skills of the 2026 programme applicable to this cohort.

\textbf{Teaching format of the booklet.} Each chapter contains two MCQs of five questions. The annual papers use twelve questions worth 0.5 point each, followed by three exercises totalling 14 points: this is a preparation choice compatible with the regulatory structure, not a claim about the number of questions in the future official paper. The June 2026 metropolitan paper, for example, has eight fluency questions.

\textbf{Calendar.} The reference group has 2-hour slots on Tuesdays and Thursdays, from 9:45 to 11:45. The common test is announced for 14 November; the mock exams for 14 to 16 January and 1 to 3 April. Their precise times are not supplied. The papers offered last 2 h; this duration remains a proposal for the common assessments. The date of the EAM in Madagascar is not deduced from the metropolitan calendar.

\textbf{Official references.}\par
\href{https://www.education.gouv.fr/bo/2026/Hebdo14/MENE2602917A}{Première programme, BO of 2 April 2026}.\par
\href{https://www.education.gouv.fr/bo/2025/Hebdo24/MENE2515469N}{Definition of the exam, BO of 12 June 2025}.\par
\href{https://eduscol.education.gouv.fr/5688/epreuve-anticipee-de-mathematiques-aux-baccalaureats-general-et-technologique}{Éduscol: presentation and sample papers}.\par
\href{https://www.apmep.fr/Epreuve-anticipee-2026}{APMEP: directory of the 2026 papers, several centres}.
\tableofcontents
\mxnewpage
